\documentclass[11pt]{article}

% =============================================================================
% CSCI 5832 final report, written as a SemEval-2027 system description paper.
%
% Class requirements: LaTeX with the ACL template, 5-8 pages, submitted as PDF,
% meeting the requirements of a real SemEval submission:
%   https://semeval.github.io/paper-requirements.html
%   https://semeval.github.io/system-paper-template.html
%
% Before submitting, check that:
%   - The title follows "<Team name> at SemEval-2027 Task 2: <title>".
%   - Author names and affiliations are filled in (SemEval papers are not anonymous).
%   - The task paper is cited; swap the placeholder entry in custom.bib for the
%     official BibTeX once the organizers release it.
%   - The introduction states the data language (English) and gives the code URL.
%   - Results cover several seeds, the baselines and ablations.
%   - Every number in the paper comes from committed code and result files.
%   - The AI-assistance statement in the Acknowledgments is complete.
%   - No red [TODO] markers remain (search this file for "\todo").
%   - The main text is 5-8 pages. In ACL style, references and appendices do not
%     count towards the page limit; check how your instructor counts.
% =============================================================================

% "final" gives the camera-ready layout with author names. Do not use "review":
% it hides the authors and adds line numbers, but SemEval papers are not anonymous.
\usepackage[final]{acl}

% Standard package includes
\usepackage{times}
\usepackage{latexsym}
\usepackage[T1]{fontenc}
\usepackage[utf8]{inputenc}
\usepackage{microtype}
\usepackage{inconsolata}
\usepackage{graphicx}

% Added for this paper
\usepackage{booktabs}  % \toprule, \midrule, \bottomrule
\usepackage{amsmath}
\usepackage{amssymb}   % \mathbb
\usepackage{xcolor}

% Visible placeholder: every \todo{...} prints in red in the PDF.
\newcommand{\todo}[1]{\textcolor{red}{[TODO: #1]}}

% Label abbreviations used in tables and text.
\newcommand{\FWD}{\textsc{fwd}}
\newcommand{\BWD}{\textsc{bwd}}
\newcommand{\EQ}{\textsc{eq}}
\newcommand{\NEG}{\textsc{neg}}

% Replace TeamName with your team's name, and use a title that says what is
% distinctive about your approach.
\title{TeamName at SemEval-2027 Task 2: Descriptive Title of Your System}

% Add co-authors with \And (different institution) or \and (same institution).
\author{Victor Jimenez \\
  University of Colorado Boulder \\
  \texttt{viji2154@colorado.edu} \\\And
  Sepehr Rezaei \\
  University of Colorado Boulder \\
  \texttt{sere8147@colorado.edu} \\}

\begin{document}
\maketitle

\begin{abstract}
\todo{Four to six sentences: the task, your approach, your main dev results (weighted F1, SoftCons, HardCons), and the main finding of your analysis.}
\end{abstract}

% -----------------------------------------------------------------------------
\section{Introduction}
% SemEval guidance: about one paragraph per point. Cite the task paper here.

\todo{The task and why it matters: systems label an ordered pair of English phrases and must give compatible answers when the pair is swapped \citep{dico-nli-2027-task}, a reversal stress test inspired by the Reversal Curse \citep{berglund2024reversal}.}

\todo{Your main strategy, in a few sentences.}

\todo{Your key results, how they compare with the baselines, and what the analysis shows.}

Our code is available at \url{https://github.com/sepehrRez/NLP-Semeval-2027}.
\todo{Make sure the repository is public before submitting.}

% -----------------------------------------------------------------------------
\section{Background}

\subsection{Task and Data}
% SemEval: describe the input, the output, the data (language, genre, size)
% and the track you entered.

\todo{Describe the task in your own words: the input, the four labels (Table~\ref{tab:labels}), where the phrases come from (image captions and news headlines from PhrasIS; \citealp{lopezgazpio2024phrasis}), the track (Track~1, English), and the data sizes (Table~\ref{tab:data}).}

\begin{table}[t]
  \centering
  \small
  \begin{tabular}{@{}lll@{}}
    \toprule
    \textbf{Label} & \textbf{Rev} & \textbf{Example (text1 $\rightarrow$ text2)} \\
    \midrule
    \FWD & \BWD & \textit{ballistic missiles} $\rightarrow$ \textit{missiles} \\
    \BWD & \FWD & \textit{missiles} $\rightarrow$ \textit{ballistic missiles} \\
    \EQ  & \EQ  & \textit{flying} $\rightarrow$ \textit{in flight} \\
    \NEG & --   & \textit{a dog} $\rightarrow$ \textit{the skateboarder} \\
    \bottomrule
  \end{tabular}
  \caption{Labels, their reversal (Rev) when text1 and text2 are swapped, and examples from the training data. \FWD{} = forward entailment, \BWD{} = backward entailment, \EQ{} = equivalence, \NEG{} = negative/other. \NEG{} items are excluded from the consistency scores.}
  \label{tab:labels}
\end{table}

\begin{table}[t]
  \centering
  \small
  \begin{tabular}{@{}lrrrrrr@{}}
    \toprule
    \textbf{Split} & \FWD & \BWD & \EQ & \NEG & \textbf{Total} & \textbf{Pairs} \\
    \midrule
    Train & 881 & 881 & 802 & 478 & 3{,}042 & 1{,}282 \\
    Dev   & 190 & 190 & 174 & 106 & 660     & 277 \\
    \bottomrule
  \end{tabular}
  \caption{Track~1 instances per label. Pairs: reversible source pairs, which appear in both directions and count towards SoftCons and HardCons.}
  \label{tab:data}
\end{table}

\subsection{Evaluation}

Let $f(x)$ be a system's prediction for instance $x$, $y$ its gold label, $x^{r}$ the same pair in the opposite direction, and $\mathrm{Rev}$ the reversal map in Table~\ref{tab:labels}. With sums over the $N$ reversible pairs $(x, x^{r})$ \citep{dico-nli-2027-task}:
\begin{align}
  \mathrm{SoftCons} &= \frac{1}{N} \sum \mathbb{I}\big[ f(x) = \mathrm{Rev}(f(x^{r})) \big] \\
  \mathrm{HardCons} &= \frac{1}{N} \sum \mathbb{I}\big[ f(x) = y \wedge f(x^{r}) = y^{r} \big]
\end{align}

\todo{Explain in your own words what each score captures, and that weighted F1 is computed over all four labels.}

\subsection{Related Work}
% Rubric: "a section created with decent effort summarizing prior work".
% Group the work by theme and say how yours relates to each. All of these are in
% custom.bib (the first group of entries comes from the organizers):
%   Task lineage: \citet{lopezgazpio2017interpretable}, \citet{lopezgazpio2024phrasis},
%                 the DiCo-NLI pilot \citet{apaolaza2026assessing}
%   NLI:          \citet{maccartney2009natural}, \citet{bowman-etal-2015-large},
%                 \citet{williams-etal-2018-broad}
%   Shortcuts:    \citet{gururangan2018annotation}
%   Reversal and consistency: \citet{berglund2024reversal}, \citet{golovneva2024reverse},
%                 \citet{lin2024delving}, \citet{jang2022becel},
%                 \citet{wu2025transitive}, \citet{blanck2026reverse}
%   Evaluation:   \citet{lopez2024revisiting}

\todo{Two to four paragraphs on prior work.}

% -----------------------------------------------------------------------------
\section{System Overview}
% SemEval: usually the longest section. Give enough detail that a reader could
% re-implement your system.

\subsection{Fine-tuned Encoders}
\todo{The baseline from the organizers' starter kit: a pretrained encoder reads text1 and text2 as one input, and a classification head on the first token predicts one of the four labels. Name the encoders you compare, e.g.\ BERT \citep{devlin-etal-2019-bert}, RoBERTa \citep{liu2019roberta}, DeBERTaV3 \citep{he2023debertav3} and ModernBERT \citep{warner2024modernbert}.}

\subsection{Consistency-Aware Inference}
\todo{Your own method, with equations, and why it should help the consistency scores.}

\subsection{Additional Systems}
\todo{Optional, e.g.\ an LLM prompted zero- or few-shot. The full prompt goes in Appendix~\ref{app:prompts}.}

% -----------------------------------------------------------------------------
\section{Experimental Setup}
% SemEval: data splits, preprocessing, hyperparameter tuning, external tools
% (versions and URLs in a footnote), and evaluation measures.

\todo{Data splits: hold out part of train as a validation set, grouped by \texttt{pair\_id} so both directions of a pair stay together. Tune and pick epochs on it, and evaluate on the official dev set only at the end. The official test set is not released until January 2027.}

\todo{Training details: optimizer, learning-rate schedule, batch size, epochs, and number of seeds (report mean $\pm$ standard deviation). Full values in Appendix~\ref{app:hyperparams}.}

\todo{Hardware and software: Apple M4 GPU through PyTorch MPS, PyTorch \citep{paszke2019pytorch}, and Hugging Face Transformers \citep{wolf-etal-2020-transformers}, with versions and URLs in a footnote.}

\todo{The shared-task rules you followed: no PhrasIS data, no training on the trial data (it overlaps with dev), and no use of the data-file structure as a feature.}

% -----------------------------------------------------------------------------
\section{Results}

\todo{Describe the main results (Table~\ref{tab:main}) and the ablations (Table~\ref{tab:ablation}): which models do best, and what the consistency scores show beyond F1.}

% Full page width (table*) leaves room for "mean ± std" in every cell.
\begin{table*}[t]
  \centering
  \small
  \begin{tabular}{@{}lccc@{}}
    \toprule
    \textbf{System} & \textbf{W-F1} & \textbf{SoftCons} & \textbf{HardCons} \\
    \midrule
    \multicolumn{4}{@{}l}{\textit{Baselines}} \\
    Majority class                 & -- & -- & -- \\
    Longer phrase entails shorter  & -- & -- & -- \\
    \midrule
    \multicolumn{4}{@{}l}{\textit{Fine-tuned encoders}} \\
    BERT-base        & -- & -- & -- \\
    RoBERTa-base     & -- & -- & -- \\
    DeBERTa-v3-base  & -- & -- & -- \\
    ModernBERT-base  & -- & -- & -- \\
    \midrule
    \multicolumn{4}{@{}l}{\textit{Our system}} \\
    \todo{name}      & -- & -- & -- \\
    \bottomrule
  \end{tabular}
  \caption{Track~1 dev results, mean $\pm$ standard deviation over \todo{N} seeds. W-F1: weighted F1.}
  \label{tab:main}
\end{table*}

\begin{table}[t]
  \centering
  \small
  \setlength{\tabcolsep}{4pt}
  \begin{tabular}{@{}lccc@{}}
    \toprule
    \textbf{Variant} & \textbf{W-F1} & \textbf{SoftCons} & \textbf{HardCons} \\
    \midrule
    Full system                       & -- & -- & -- \\
    \quad $-$ \todo{component}        & -- & -- & -- \\
    \quad $-$ \todo{component}        & -- & -- & -- \\
    \bottomrule
  \end{tabular}
  \caption{Ablations on Track~1 dev. Each row removes one component from the full system.}
  \label{tab:ablation}
\end{table}

% -----------------------------------------------------------------------------
\section{Analysis}
% Rubric: 20 points for analysing the results (error analysis, follow-up
% experiments). SemEval: "analysis of your system's behavior beyond the overall scores".

\subsection{Error Analysis}
\todo{Per-label scores and the confusion matrix (Figure~\ref{fig:confusion}): which labels get confused, and how often the direction is wrong.}

\begin{figure}[t]
  \centering
  \includegraphics[width=\columnwidth]{example-image-golden}
  \caption{\todo{Confusion matrix of your best system on dev; replace the placeholder image.}}
  \label{fig:confusion}
\end{figure}

\subsection{Consistency Failures}
\todo{Which pairs get incompatible answers in the two directions, and why.}

\subsection{Shortcuts}
\todo{How often each system agrees with the ``longer phrase entails the shorter one'' heuristic, and where it does better (e.g.\ numbers, names, paraphrases).}

\subsection{Manual Inspection}
\todo{Categorize about 50 errors by hand and report the categories, with examples.}

% -----------------------------------------------------------------------------
\section{Conclusion}
\todo{A few sentences: what you built, what you found, and future work.}

% Standard in ACL papers; it does not count towards the page limit.
\section*{Limitations}
\todo{E.g.\ English only, results on the dev set rather than the official test set, a single hardware setup, and the size of the training data.}

% ACL asks authors to disclose AI assistance beyond minor polishing. Keep this
% statement complete and accurate as the project continues.
\section*{Acknowledgments}
We used Claude Opus 5.5 (Anthropic), an AI assistant, through Claude Code to set up the software environment, adapt the organizers' starter kit to run on an Apple-Silicon GPU, draft a project plan, and create the \LaTeX{} skeleton of this paper.
\todo{Add any later AI use (code, analysis, writing), and thank anyone who helped.}

\bibliography{custom}

\appendix

\section{Hyperparameters}
\label{app:hyperparams}

\todo{Search space and selected values for each model (Table~\ref{tab:hyperparams}).}

\begin{table}[h]
  \centering
  \small
  \begin{tabular}{@{}lll@{}}
    \toprule
    \textbf{Hyperparameter} & \textbf{Search space} & \textbf{Selected} \\
    \midrule
    Learning rate          & -- & -- \\
    Batch size             & -- & -- \\
    Epochs                 & -- & -- \\
    Warmup ratio           & -- & -- \\
    Weight decay           & -- & -- \\
    Dropout                & -- & -- \\
    Max.\ sequence length  & -- & -- \\
    Seeds                  & -- & -- \\
    \bottomrule
  \end{tabular}
  \caption{Hyperparameters for \todo{model}.}
  \label{tab:hyperparams}
\end{table}

\section{Prompts}
\label{app:prompts}

\todo{Only if you use an LLM: the exact prompt(s) and decoding settings.}

\end{document}
